\section{0. Program Thesis, Claim Boundary, and Implementation Contract}

This document is not a report that SC-WBD has already achieved whole-brain
prediction. It is the technical thesis that governs what SuperCognition Labs
will build, what evidence may update it, which configurations the compiler
must reject, and which experiments are capable of changing the program's
scientific commitments. Its first audience is therefore dual: researchers
should be able to locate the assumptions that require evidence, while
implementation agents should be able to translate each contract into schemas,
validators, tests, and acceptance gates.

The thesis has four load-bearing commitments. First, heterogeneous methods are
related through a latent generative process and method-specific read/write
operators, not through naive resampling into a common tensor. Second, adult
human anatomy supplies strong probabilistic topology and phenotype priors,
while effective coupling, local computation, state, and slow plasticity remain
learnable. Third, resolution is an indexed property of state and evidence:
several non-equivalent spatial and temporal supports may coexist without a
single materialized truth. Fourth, mechanistic language is earned only by
predictions that an equal-capacity surrogate misses, especially under held-out
perturbation. These commitments are compatible with useful predictive models
that remain agnostic about many biological details; they do not license a
claim that every admitted operator is neurally realized.

\begin{table*}[t]
\centering
\footnotesize
\setlength{\tabcolsep}{4pt}
\renewcommand{\arraystretch}{1.10}
\begin{tabularx}{\textwidth}{@{}p{.19\textwidth}p{.27\textwidth}p{.27\textwidth}X@{}}
\toprule
\textbf{Program claim} & \textbf{Evidence that can support it} &
\textbf{Result that weakens or falsifies it} & \textbf{Implementation consequence}\\
\midrule
Typed, source-native fusion is preferable to naive resampling &
Better held-out likelihood, calibration, delay recovery, and intervention
forecast at matched compute and parameter count, measured against a tuned
resampling baseline. Adding a modality is not itself evidence: under (T4) it
cannot reduce information &
No reproducible gain over single-modality or carefully tuned resampling
baselines; increased overconfidence or negative transfer &
Retain only the provenance/type system; narrow or remove the shared latent
fusion claim. \emph{Narrowed on measured evidence to a claim about
source-native clocks; see Section 11.5}\\

Anatomical topology improves inference &
Data efficiency, out-of-distribution calibration, and causal forecast beyond
dense, randomized, and distance-matched graphs &
Equal-capacity controls match or exceed performance, or topology errors are
absorbed by residuals &
Demote anatomy from compiled constraint to weak prior for the affected scale\\

Multiresolution state adds information rather than decoration &
Native-scale prediction, calibrated refinement, boundary agreement, and
compute savings without hidden high-frequency hallucination &
Fine views cannot improve supported observables, fail round-trip tests, or
become overconfident outside measured tiles &
Disable the scale relation or use source-specific views without gluing\\

Perturbation reduces non-identifiability &
Rank/eigenvalue improvement in Fisher information \emph{at matched input
energy}, and prospective recovery of direction, delay, gain, dose, and state
dependence &
Intervention fails to distinguish posterior models, adds only field-model
uncertainty, or loses its apparent advantage against an energy-matched control &
Narrow the identifiable parameter set and redesign the perturbation rather
than reporting a causal estimate\\

Individualization improves future prediction &
Incremental calibrated log score and decision utility on new sessions and
unseen tasks or interventions &
Anatomy-only, population, or session-adapted baselines perform equivalently &
Do not label the model an individual twin; retain only the supported level of
adaptation\\
\bottomrule
\end{tabularx}
\caption{\textbf{Claims with consequences.} Each central commitment has a
measurement capable of making it smaller. Engineering breadth, parameter
count, plausible diagrams, and in-sample fit are not substitutes for these
tests. Rows 1 and 4 have been narrowed since the previous version because the
measurement was run and returned a smaller result than the claim; the evidence
columns record the controls that narrowing showed to be mandatory
(Section 11.5).}
\label{tab:claim-gates}
\end{table*}

\subsection{0.1 What SC-WBD Forbids}

A type system earns credibility by rejecting programs. SC-WBD therefore
distinguishes an operator family that can be represented from a model instance
that is admissible for a particular scientific claim. The compiler must fail
closed on the following conditions unless the schema explicitly changes the
claim class and records the override.

\begin{table*}[t]
\centering
\footnotesize
\setlength{\tabcolsep}{4pt}
\renewcommand{\arraystretch}{1.09}
\begin{tabularx}{\textwidth}{@{}p{.22\textwidth}p{.31\textwidth}X@{}}
\toprule
\textbf{Rejected configuration} & \textbf{Why it is scientifically invalid} &
\textbf{Required remedy}\\
\midrule
Unknown units, clock, physical support, frame, handedness, or transform lineage &
Numerical compatibility would be mistaken for measurement compatibility &
Supply a calibrated manifest with validity interval and uncertainty, or keep
the source disconnected from the requested path\\

Prolongation without a declared restriction partner and tested coverage &
A coarse observation would be converted into unsupported fine structure &
Provide paired maps, round-trip and held-out landmark tests, and an
out-of-support uncertainty policy; otherwise return a distribution or refuse\\

Global cross-scale state when overlap/cocycle residual exceeds tolerance &
Conflicting local views would be hidden by one raster &
Preserve separate sections, branch or marginalize the posterior, and emit an
obstruction certificate identifying the failed path\\

Effective or causal operator estimated from passive correlation alone &
Common input, observation mixing, and omitted state remain alternatives &
Label the object functional/statistical, or add an identified intervention or
assumption set sufficient for the causal claim\\

Learned residual allowed to dominate a mechanistic term silently &
The named mechanism becomes unfalsifiable &
Enforce a preregistered energy/gain ratio
\(\|\mathcal R_\theta\|/\|\mathcal F_{\rm mech}\|\leq\rho_{\max}\) on the
declared validity set, report violations, or reclassify the entire module as a
surrogate\\

Adaptive step sizes for a learned propagator without semigroup testing &
Coarse and composed fine transitions need not describe the same dynamics &
Measure the semigroup residual at every permitted step pair; restrict the
scheduler or include the discrepancy in prediction intervals\\

Population, subject, and session effects without centering or shrinkage &
The additive decomposition is not identified &
Use centered/sum-to-zero constraints or proper hierarchical priors and test
recovery in simulation\\

Bias assigned a point estimate without an estimator or external bound &
Calibration and model discrepancy can trade off non-identifiably &
Classify the term as design-estimable, externally bounded, or
prior-specified sensitivity; propagate the last category as a range\\

Pseudo-likelihood compatibility treated as a calibrated posterior likelihood &
Agreement penalties do not define a generative noise model &
Report a generalized/pseudo-posterior and calibrate it empirically, or validate
and promote the factor to a generative likelihood\\

Derived scans, sessions, relatives, stimuli, or simulator replicas crossing a
parent-level holdout &
Duplicated evidence produces invalid confidence and shortcut prediction &
Group by immutable lineage identifiers before splitting and fail the run when
parentage is unresolved\\

Intervention optimization outside an independently validated feasible set &
A learned objective cannot guarantee device, exposure, or protocol safety &
Require \(u\in\mathcal A_{\rm safe}\), independent safety checks, deferral,
and the applicable ethics and regulatory review\\
\bottomrule
\end{tabularx}
\caption{\textbf{Mandatory modeling refusals.} These are compiler and runtime
errors, not optional governance prose. An override changes the claim carried
by the resulting artifact and remains visible in its provenance.}
\label{tab:compiler-refusals}
\end{table*}

\subsection{0.2 Position Relative to Existing Model Families}

The latent-state/observation-operator decomposition is not new. Dynamic causal
modeling (DCM) already formalizes Bayesian inference over hidden neuronal
dynamics, effective connectivity, experimental inputs, and modality-specific
forward models \citep{friston2003dcm}. The Virtual Brain provides an
individualizable connectome-coupled simulation environment
\citep{sanz2013,sanz2015}; neurolib and BrainPy provide extensible machinery
for whole-brain neural-mass and general brain-dynamics programming
\citep{cakan2023neurolib,wang2023brainpy}. NeuroML/LEMS and PyNN establish
substantial prior art for declarative, unit-aware, simulator-independent model
specification \citep{cannon2014lems,davison2009pynn}. SC-WBD should interoperate
with these ecosystems rather than redescribe them as absent.

The proposed differentiator is the joint compilation of (i) heterogeneous
operator-valued regional state; (ii) non-nested, source-native spatial and
temporal resolutions; (iii) uncertain coordinate/clock/field transforms;
(iv) data-use provenance into gradient permissions and leakage barriers; and
(v) explicit claim-dependent refusal rules. DCM remains the clearest
statistical ancestor for generative inversion; declarative neuroscience
languages are the clearest software ancestor for typed execution. SC-WBD's
burden is to demonstrate that extending those ideas to multirate,
multiresolution, intervention-aware fusion yields identifiable and calibrated
predictions unavailable to simpler constructions.

Patient-specific cardiac modeling is a useful neighboring discipline because
it has had to separate code verification, physical validation, parameter
calibration, model-form uncertainty, and context-of-use credibility before a
digital-twin claim can influence a decision \citep{galappaththige2022}. SC-WBD
adopts the same ladder: numerical correctness is necessary; agreement with
recorded signals is stronger; held-out perturbation is stronger still; and
prospective decision improvement is a distinct claim.

\subsection{0.3 Identifiability Before Scale}

The first scientific artifact is intentionally small. Consider three latent
regional states with directed coupling, one unknown conduction delay, and a
controlled input:

\[
\begin{aligned}
\mathrm d x(t)={}&A(\vartheta)x(t)\,\mathrm dt
+B\,u(t-\tau_u)\,\mathrm dt\\
&+Q^{1/2}\mathrm dW_t,
\qquad x(t)\in\mathbb R^3 .
\end{aligned}
\tag{T1}
\]

An EEG-like instrument observes rapid instantaneous mixing,

\[
y_E[k]=L(\ell)x(k\Delta_E)+\epsilon_E[k],
\qquad \Delta_E=1~\mathrm{ms},
\tag{T2}
\]

whereas an fMRI-like instrument observes a slow hemodynamic convolution,

\[
y_B[n]=M\!\int_0^{T_h}h(s;\rho)x(n\Delta_B-s)\,\mathrm ds
+\epsilon_B[n],
\qquad \Delta_B=1~\mathrm{s}.
\tag{T3}
\]

The unknown vector \(\vartheta\) contains selected coupling gains and the
network delay, while \(\ell\) and \(\rho\) contain observation nuisance
parameters. For a linear--Gaussian discretization the expected Fisher
information for design \(d\) is

\[
\mathcal I_d(\eta)=
\sum_{m\in d}J_m(\eta)^\top R_m^{-1}J_m(\eta)
+\mathcal I_{\rm prior},
\quad
J_m=\frac{\partial\mu_m}{\partial\eta},
\tag{T4}
\]

with \(\eta=(\vartheta,\ell,\rho)\).

Two properties of (T4) must be stated where it is introduced, because a
benchmark that ignores either of them will report an artifact of its own
construction as a finding.

\textbf{Modality additivity is an identity, not a hypothesis.} Written as a sum
of per-modality quadratic forms, (T4) treats the modalities as conditionally
independent given \(\eta\). Each summand is positive semidefinite, so
\(\mathcal I_{\rm EEG+BOLD}
=\mathcal I_{\rm EEG}+\mathcal I_{\rm BOLD}-\mathcal I_{\rm prior}
\succeq\mathcal I_{\rm EEG}\)
holds algebraically, for every parameterization, every dataset, and every
truth. ``Joint is at least as informative as single-modality'' therefore cannot
fail, and a benchmark reporting it as a validated result reports a tautology.
The non-trivial part of joint information is the EEG/BOLD cross-covariance
induced by \emph{shared process noise}, which the block-diagonal form omits;
obtaining it requires whitening the stacked record rather than each modality
separately. The benchmark computes both and reports their difference, which is
exactly the amount by which (T4) over-counts. The discriminating comparisons
are consequently native versus resampled clocks, the \emph{size} of the second
modality's contribution to the preregistered subset, impulse versus no impulse
at matched input energy, and whether any added information becomes calibrated
recovery.

\textbf{A continuous conduction delay is a precondition for delay
identifiability.} The Fisher information of \(\tau\) exists only if
\(\partial\mu/\partial\tau\) exists. If \(\tau\) is quantized to the sample
grid, that derivative is identically zero and the delay carries \emph{no}
information under any design---including designs that would otherwise recover
it well. A benchmark built on a quantized delay would report zero delay
information for every design and mistake an implementation choice for a
scientific result. (T1) and (T4) therefore require a fractional-delay
implementation; here a normalized Gaussian-windowed-sinc interpolation over the
history taps, so that \(\tau\) is a continuous parameter with a smooth
derivative. The 1\,s resampled design is retained precisely because it violates
this condition, and is reported as the illustration of what the violation looks
like rather than as a measurement of the delay.

The benchmark compares EEG alone, fMRI alone, joint native-clock inference,
joint inference after naive resampling, joint inference plus one calibrated
impulse, and---required rather than optional---the same impulse at total input
energy matched to the impulse-free design by scaling the background drive.
Without that sixth design, an impulse comparison measures the energy budget and
not the perturbation. A charitable resampling control, which discards EEG onto
the 1\,s grid but retains the exact fine-grained forward model, separates
information lost by decimation from bias introduced by the wrong
discretization. The benchmark reports rank, condition number, minimum non-prior
eigenvalue, parameter-profile likelihoods, posterior correlations, interval
coverage, and delay error. Prior contribution is shown separately so a
full-rank posterior cannot disguise a prior-dominated likelihood.

\textbf{What this system cannot be asked.} The state is three regions with no
spatial structure, and \(\vartheta\) is a few coupling gains plus one
conduction delay. That is an electrophysiological question, and a hemodynamic
instrument is being scored on it. The strengths such an instrument actually
has---spatial extent and localization, slow state, and constraining the very
observation model on which electrical source inversion depends---have nowhere
to appear in a three-region linear--Gaussian system. A weak fusion result
obtained here is therefore a correct statement about \emph{this design}. It is
not a general finding that cross-method integration is unprofitable, and it may
not be reported as one; reaching that stronger question requires extending the
system, not rerunning this one.

This is a planned calculation, not a favorable result assumed in advance. The
central premise is supported only if native-clock fusion or intervention
increases likelihood information for a preregistered parameter subset and
improves calibrated recovery across held-out simulation regimes. If it does
not, the compiler may still be useful as a provenance system, but the claim
that cross-method integration resolves those dynamics must be narrowed. That
calculation has now been performed in the linear--Gaussian case; its outcome,
and the narrowing this paragraph prescribes, are reported in Section 11.5. The
paragraph is left exactly as written, because it is the commitment against
which the result is being judged.

The second artifact compiles the same system from a source schema and performs
end-to-end recovery. It introduces known transform bias, shared session
calibration error, missing windows, unequal supports, and a deliberately
misspecified residual. Success requires: no leakage across simulated parent
subjects; recovery intervals with nominal coverage; lower held-out log loss
than single-method and naive-resampling baselines; detection of the
misspecified module; and refusal of at least one invalid schema. This synthetic
case is the minimum integration test for all later human experiments. It has
been built and run: three of the five criteria listed here were met, interval
coverage was not, the held-out log-loss comparison could not be evaluated
because none of the fits it compares had converged, and a sixth
subgroup-calibration criterion added by the implementation was met
(Section 11.5).

\subsection{0.4 Executable Evidence, Transform, and Gradient Contracts}

Every dataset snapshot is compiled from a source card rather than passed to a
global dataloader. A source card declares identity and hashes, participant and
derivative lineage, governance, native signal and units, spatial support,
point-spread function, clock, coordinate/calibration graph, observation or
intervention model, missingness, validity domain, bias/variance/discrepancy
ledger, permitted module and scale targets, and a fixed role: prior,
likelihood, boundary target, distillation, calibration, negative control, or
evaluation only. The compiler converts those declarations into dataloaders,
transform paths, loss factors, gradient masks, and group-aware splits. An
unknown field remains unknown; it is not supplied as zero or as an average
brain label.

The uncertainty ledger uses three different statuses. A quantity is
\textbf{design-estimable} only when the acquisition contains the replication,
randomization, intervention, anchor, or hierarchy necessary to estimate it.
Repeated sessions, for example, can estimate within- and between-session
variance but do not by themselves identify absolute model bias. A term is
\textbf{externally bounded} when a phantom, calibration target, independent
instrument, or negative control supplies a defensible interval. A term is
\textbf{prior-specified sensitivity} when neither is available; it is swept
over a declared range and cannot be advertised as empirically estimated.
Model discrepancy and calibration parameters are generally confounded without
informative assumptions, a known failure mode in computer-model calibration
\citep{kennedy2001,brynjarsdottir2014}.

For a derived quantity \(z=f(x,c)\), calibration uncertainty may be shared
across every observation in a session. First-order propagation therefore
retains cross-covariance:

\[
\begin{aligned}
\Sigma_z\approx{}&J_x\Sigma_xJ_x^\top+J_c\Sigma_cJ_c^\top\\
&+J_x\Sigma_{xc}J_c^\top+J_c\Sigma_{cx}J_x^\top,\\
b_z\approx{}&J_xb_x+J_cb_c+\delta_f .
\end{aligned}
\tag{T5}
\]

The transform graph stores each edge separately. Paths are checked for units,
handedness, epoch, support, round-trip residual, and calibration validity;
their composed result never replaces the original edges. Figure
\ref{fig:framegraph} gives the frame path used by the running example.

\FrameGraphFigure

For episode \(e\), only authorized additive factors are active:

\[
\begin{aligned}
\mathcal L=\sum_e w_e\bigl[&
\lambda_{\rm like}\mathcal L_{\rm like}^{(e)}
+\lambda_{\rm boundary}\mathcal L_{\rm boundary}^{(e)}\\
&+\lambda_{\rm distill}\mathcal L_{\rm distill}^{(e)}
+\lambda_{\rm cal}\mathcal L_{\rm cal}^{(e)}\\
&+\lambda_{\rm mech}\mathcal R_{\rm mech}^{(e)}
+\lambda_{\rm scale}\mathcal R_{\rm scale}^{(e)}\bigr].
\end{aligned}
\tag{T6}
\]

Weights depend on effective independent evidence and declared reliability, not
row count. Gradient conflict is logged by source and module. A conflict may
produce partial pooling, a source-specific adapter, a frozen path, a branched
posterior, or rejection; consensus is not mandatory.

\subsection{0.5 Running Case: Individual DLPFC TMS Targeting}

The running case asks a bounded question: for one consented adult in a
specified session and cognitive state, which of a small preregistered set of
left-DLPFC coil poses is predicted to produce the desired short-horizon change
in a measured network response while remaining inside independent exposure
and protocol limits? It does not ask whether the model treats depression, nor
does it infer a latent clinical construct from one signal.

\begin{enumerate}
\item \textbf{Declare.} The schema selects the subject's cortical surface,
  uncertain tract topology, local DLPFC field, coarse distributed network,
  EEG and fMRI observation heads, coil geometry, neuronavigation frames,
  stimulation waveform, session state, and \(\mathcal A_{\rm safe}\).
\item \textbf{Compile.} The compiler checks source cards and transform paths,
  instantiates a fine DLPFC tile plus coarse surrounding regions, registers
  restriction/prolongation maps, and emits native EEG and fMRI clocks. A pose
  expressed only as ``5 cm anterior'' is rejected.
\item \textbf{Calibrate.} MRI constrains geometry; digitized electrodes and
  head model constrain the EEG lead field; tracker-to-coil and head-to-MRI
  observations constrain pose; motor-threshold or other approved calibration
  informs device gain without being equated to DLPFC target engagement.
\item \textbf{Infer.} Baseline recordings produce a posterior over neural
  state, coupling, observation nuisance, tissue parameters, and model class.
  Session calibration errors remain correlated rather than being averaged
  away. Models that fit passive data but disagree about propagation are
  retained.
\item \textbf{Counterfactually compare.} Each admissible pose is passed through
  the electromagnetic field solver and candidate neural response operators.
  The system predicts a distribution over rapid EEG propagation, later
  hemodynamic consequences, burden, and explicit harm endpoints. Field
  accuracy, target engagement, network effect, and clinical utility remain
  separate quantities.
\item \textbf{Decide or defer.} The optimizer acts only within
  \(\mathcal A_{\rm safe}\). If model disagreement or transform uncertainty
  dominates the estimated benefit difference, the output is another
  calibration measurement, a reversible probe, or no recommendation.
\item \textbf{Test.} A prospective, blinded or otherwise causally identified
  comparison evaluates predicted direction, latency, spatial propagation,
  dose response, state dependence, calibration, and decision regret. The
  result updates only the parameters and claims authorized by the protocol.
\end{enumerate}

Prospective human TMS or tFUS work requires ethics, safety, consent, and device
review appropriate to the actual protocol. In the United States the sponsor
and IRB must document significant-risk versus nonsignificant-risk status; a
significant-risk device investigation requires FDA and IRB approval, while an
NSR study follows the abbreviated IDE requirements after IRB concurrence
\citep{fdaide}. IRB approval alone is therefore not treated as a universal
regulatory description.

\subsection{0.6 Agent-Facing Build Order and Definition of Done}

Implementation proceeds through claim-bearing vertical slices rather than a
simultaneous attempt at every brain system.

\begin{enumerate}
\item \textbf{Schema kernel.} Implement typed regions, ports, operators,
  clocks, supports, frames, authority policies, uncertainty status, source
  roles, and immutable lineage. Definition of done: valid examples compile;
  every refusal in Table~\ref{tab:compiler-refusals} has a failing fixture.
\item \textbf{Linear identifiability laboratory.} Implement Equations
  (T1)--(T4), analytic/numerical Fisher information, simulation, posterior
  recovery, and five design comparisons. Definition of done: deterministic
  seeds, parameter sweeps, coverage tests, and a machine-readable claim report.
\item \textbf{Transform and clock runtime.} Implement frame-graph composition,
  cross-covariance, multirate events, validity intervals, and obstruction
  certificates. Definition of done: unit/handedness/clock errors are caught;
  round trips and shared calibration error are tested.
\item \textbf{End-to-end synthetic slice.} Compile three regions, EEG/fMRI
  reads, one impulse write, missing windows, and model discrepancy. Definition
  of done: the preregistered baselines and acceptance criteria in Section 0.3
  run from one manifest.
\item \textbf{Empirical subsystem.} Add one anatomically bounded human
  subsystem with public data, source-native likelihoods, and held-out
  cross-modal prediction. Definition of done: provenance, leakage, bias status,
  and negative-transfer ablations ship with the result.
\item \textbf{Prospective perturbation pilot.} Only after solver, field,
  calibration, and causal recovery gates pass, instantiate the running TMS or
  tFUS case under an approved protocol. Definition of done: preregistered
  prospective predictions and an explicit no-benefit/no-identification outcome
  path.
\end{enumerate}

Every implementation agent receives the thesis version, schema version,
target claim, allowed source cards, acceptance tests, and non-goals. A merged
component must add tests, units and frame declarations, uncertainty behavior,
provenance fields, baseline comparisons, and a statement of what empirical
finding would disable it. Code that only expands the operator registry is not
scientific progress until it passes a claim-bearing gate.
